\section{Evola and Nietzsche – 40 years later}

\begin{quotex}
Transcending oneself: this is the great imperative of the human condition; and there is another that anticipates it and at the same time prolongs it: dominating oneself. The noble man is the one who dominates himself; the holy man is the one who transcends himself.
\flright{\textsc{Frithjof Schuon}}
\end{quotex}


In writing about the significance of \textbf{Friedrich Nietzsche} today, it is clear that \textbf{Julius Evola} is more concerned about the significance of Nietzsche for Evola. This is actually an important question, since it is through Evola that Nietzsche has entered into Traditional thought, at least in some circles. Usually this has had a deforming effect because Nietzsche's ideas trump Tradition; Nietzsche needs to be adapted to Tradition, not the reverse. This is what Evola tries to do. For Evola, the ``figure of Nietzsche'' is more important as a ``symbol'' rather than for any specific doctrine.

First of all, any of Nietzsche's metaphysical claims need to be rejected, and Evola never adopts them. Specifically, eternal return is a doctrine that is incompatible with the Infinite; Guenon explicitly rejected it. Obviously, Nietzsche's naturalism is not part of any Tradition, and Evola often rejected that, especially in any biological or naturalistic interpretations of the ``superman''.

Of course, Nietzsche's atheism is not helpful, which is related to the Will to Power. Together, they amount to the rejection of Intelligence and a Cosmic Order as the ultimate reality of the world. The world, then, is absurd, the result of the will to power. In the Traditional view, the world, too, is absurd, since it is the result of an illusion or a fall. The task, therefore, is self-transcendence, to overcome the world. Yet, Nietzsche's naturalism does not recognize any such transcendence; hence, the world can only be overcome by more power. Unfortunately, that is a Sisyphean task and can only lead to insanity. So if Nietzsche is not a metaphysician, that just leaves Nietzsche as Psychologist, Poet and Moralist.

There is no doubt of Nietzsche's power as a psychologist, in his ability to unearth the hidden motivations of the human mind and to expose hypocrisies; whence his insightful critiques of the bourgeois society of his era. However, Nietzsche's message gets distorted in the minds of the less skilled thinkers, who reason along these lines: ``The bourgeois Christian morality is the mentality of the herd; I am a pagan; hence I don't practice the herd morality.'' This is to confuse outward labels with interiority, and to accept the false idea that joining a different tribe will \emph{ipso facto} make one a different man. Not at all, the task of self-transcendence and self-domination still remains.

Even if for no other reason, Nietzsche's skill as a poet will keep his memory alive. He writes in a powerful, aphoristic style that tries to reach the deeper and nobler parts of our emotional life. Evola often imitates that, which accounts for his enduring appeal. Such evocations to our noble and warrior instincts are quite effective with those who would likely belong to the Kshatriya caste, although less effective on Brahmins.

That leaves Nietzsche as moralist, and this is precisely the area that Evola focuses on. Here, I'm afraid, Evola engages in a lot of wishful thinking. He even claims to understand Nietzsche better than Nietzsche understood himself. In \emph{Revolt against the modern World}, Evola writes in a footnote:

\begin{quotex}
The only modern thinker who comes close to this view, yet without being aware of it, was Nietzsche; he developed a view of absolute morality with a naturalistic basis.
\end{quotex}


What this really means is that Evola has developed a view of absolute morality, but feels it would ``sell'' better if people believed it really came from Nietzsche. Evola's view is actually very sound, despite his taking pains to hide its true source. It is true that modern thinkers do not come close to Evola's view, although ancient and medieval thinkers do. In his essay, Evola admits that ``certain connections could even be established with ancient Stoic ethics, which likewise advocated an interior sovereignty.'' Once that connection is admitted, it is easy to connect it to the ancient Greek world and then to the Thomist ethics that he grew up with.

First of all, let us look at the rudiments of Evola's morality and then see how he then ``improves'' on some of Nietzsche's more famous images, and I'll concede it is an improvement. Evola summarizes his position best in the passage from Revolt:

\begin{quotex}
Undoubtedly, there may have been a margin of indetermination even in the case of traditional man, but this margin in him only served to emphasize the positive aspect of these two sayings: ``Know yourself'' and ``Be yourself'', which implied an action of inner transformation and organization leading to the elimination of this margin of indetermination and to the integration of the self. To discover the dominating trait of one's form and caste and to will it, by transforming it into an ethical imperative and, moreover, to actualize it ritually through faithfulness in order to destroy everything that ties one to the earth (instincts, hedonistic motivations, material considerations, and so on).
\end{quotex}


For Evola, there is a margin, or gap, between knowing and willing, that man, particularly modern man but even traditional man, has not fully willed to be who he is, or in other terms, has not actualized all his possibilities. Here Evola follows Tradition in that a man is born with a certain essence, including his caste, his sex, his family, his nationality, his race, and so on. His task, then, is to know his essence, \emph{who he is}, and to conform his will to his being. This is opposite to modernity, which has embraced existentialism. For the modern mind, existence is prior to essences, which are then simply arbitrary, accidental, socially constructed, and often unjust qualities a person is stuck with. Modernity, thus, is implicitly atheistic because it denies a just social order.

However, Stoicism is indeed close to Evola's view, using \textbf{Marcus Aurelius} as its exemplar. For him, the cosmos is ordered by Intelligence, or the Logos, which he calls reason, and it is inherently just. Furthermore, everything has a ``nature'', and the moral course is to conform oneself to it. In particular, a man has a nature, which may differ from an emperor (Aurelius) to a slave (Epictetus). Each has to actualize his life in a different way. For the stoic, nature is what something should be, i.e., its essence, and the margin between what actually is, and what its nature is, must be closed through an act of the will. This agrees with Evola, but does it conform to Nietzsche?

Kant denied that essences, as noumena, could be known, and only appearances, or phenomena, are known. It seems that Nietzsche goes a step further, denying essences outright, and only appearances matter. Hence, the one thing necessary is to give ``style'' to one's character\footnote{\url{https://www.gornahoor.net/?page\_id=5545\#Needful}}. When Traditional man looks within, he sees who he is and recognizes his destiny to live out that life. In doing so, he first finds himself trapped in those base and vulgar impulses that tie him to the earth. So, he responds to Zarathustra, ``I set myself free so I can become who I am.'' Evola takes pains to show that this is not anarchy, nor the rejection of authority, as so many Nietzscheans may think ``because they do not have in themselves a higher principle that commands''. This is a significant clarification of the will to power, since Evola makes clear its actual role. The will to power, the superman, or being beyond good and evil, are meaningless concepts unless they derive from ``a higher sanction''.

There is the question of the margin or gap which, as we pointed out previously\footnote{\url{https://www.gornahoor.net/?p=5471}}, are an affront to a certain desire to live fully in nature. Here we mean ``nature'' in its contemporary meaning as whatever actually exists apart from the artifacts of man. Evola objects to the phrase, ``blond beast of prey'', which seems to imply an immediate mode of existence without any such margin. It brings to mind \textbf{William Blake}'s, ``One law for the lion and ox is oppression,'' the precise opposite of a universal morality, or any morality at all, for that matter. There is no vice in the lion for killing the ox, not any virtue in the ox for being the sacrificial victim.

Evola also claims that Nietzsche's call for ``fidelity to the earth'' is not the justification for being tied to the earth and indulging the baser motivations (instincts, hedonism, consumerism). Obviously, the meaning of Nietzsche today for many is precisely to justify those motivations; we will take Evola's side on this question.

Finally, Evola addresses Nietzsche's iconoclasm because it makes him appear to justify destruction and revolution for their own sake. Evola rejects the ``revolution of nothing'' and claims that Nietzsche is merely using rhetorical techniques to appear shocking or sensational. His real target, in Evola's view, is really ``petty morality'' and ``herd morality'', in order to make room for the higher morality of the superman. It should not be necessary to point out, however, that many Nietzscheans today simply stop at the point of idol smashing and immoralism, i.e., those who cannot recognize any higher principle within themselves. I suppose this is the ``danger'' that Evola refers to.


\flrightit{Posted on 2012-12-31 by Cologero}

\begin{center}* * *\end{center}

\begin{footnotesize}
\begin{sffamily}
\texttt{Sigurd on 2017-03-04 at 05:18 said:}

What is the correlation between essence and the Kantian ``noumena'' here? Kant denies essence is immanent, in this sense? Is the Stoa essence the same as the Aristotelian essence?

\hfill

\texttt{Matt on 2017-03-11 at 16:38 said:}

Sigurd,

In Kant's system, the noumena means essence -- the ``thing-in-itself''. Kant denied that we can know a thing's essence/intrinsic nature, or substantial form (to use aristotelian and scholastic terminology). For Kant, all we can know of a thing is it's appearance (phenomenon) to us. Here, Kant breaks with the classical-scholastic tradition because that tradition asserts that we can indeed know essences, and so therefore the world is intelligible; but the implication of Kant's position is that the world is not intelligible to us because we can never know essence (the noumena), and essence is what provides intelligibility. The world remains opaque to our understanding.

As to your second question, while I have not dived into stoic philosophy to the extent I have with Aristotelian and Platonic philosophy, from what I gather in my readings, essence in Stoic philosophy is more or less used in the same sense as seen in Aristotelian thinking.

\hfill

\texttt{Sigurd on 2017-03-15 at 05:05 said:}

Very interesting. I did not know Kant's noumena was the Aristotelian essence! And here I thought Kantian thought was more-or-less reconcilable with Traditionalist metaphysics vis-a-vis Evola's influence by the post-Kantians (Schelling, Fichte), though I realize their raison d'etre was to eliminate the noumena by bringing objectivity and, in Schelling's case, naturalism to the Kantian thought.

\hfill

\texttt{Matt on 2017-03-15 at 17:53 said:}

Sigurd,

I haven't read much of Fichte and Schelling, so I am not familiar with the details of their respective systems; just the general outlines. I think Cologero is quite familiar with them, though.

With respect to the idea of reconciling Kant with Tradition, while I don't think Kant's system of thought in total* (*emphasis on that) can be reconciled with Tradition, I think there is a possibiliy that certain elements of his thought that can be legitimately synthesized with Traditional thinking. There is a strand of Thomism known as transcendental thomism whose aim is to synthesis Thomism with certain elements of Kant's epistemology. Awhile back, Cologero posted his review of one of the main texts of its founder, Joseph Marechal. It's debatable whether or not the adherents of this strand actually accomplished this aim. Many think they failed and only created a Christianized form of Kantian philosophy. Personally, while I lean towards that assessment, I'm not closed off to the possibility that this synthesis can be done.

\hfill

\texttt{Sigurd on 2017-03-16 at 09:36 said:}

I see. I understand Guenon criticized the West's tendency to be ``systematic,'' and I take this to mean turning our ``quanitfying'' gaze toward philosophy, creating some monstrous chimera full of quasi-organic axiomatic premises, some false vector that removes any quality from our experience (and I think we can refer here to something like economics), however, I am more disposed toward Western thinking than Eastern; and despite the relative non-importance of linguistic theatricism, I do think Western philosophy is meaningful if taken in its, perhaps, Platonic sense: a life, as opposed to a classification of discussion.

Well, that was quite a drawn out preface just to ask you: what philosophical system do you think best fits the Traditional mindset? First coming to the Traditionalist was interesting for me as I hadn't read anything like it before, but I also came to read the Neoplatonists and Aristotle *through* them, so perhaps it would have been more familiar than it was, as Guenon and the others *were* writing to a Western audience. Do you think something like Thomism is totally reconcilable with Tradition? And does this make Tradition philosophically and morally realist?

Thanks for continuing this discussion by the way. I'm so glad I found this website, often it feels all my Traditional interlocutors are effectively retired! It's good to see the Tradition is still alive.


\hfill

\end{sffamily}
\end{footnotesize}

